\chapter{Conclusion}

\section{Summary}
As for the use of frequency domain in super resolution tasks, we presents WaveMixSR, a wavelet-based model for image super-resolution that operates in the frequency domain. It uses the 2D-DWT for token mixing, enabling it to reconstruct high-resolution images while capturing key image properties like scale invariance, shift invariance, and sparse edges. Unlike CNNs, which are limited in how quickly they expand their receptive fields, WaveMix efficiently mixes distant tokens. And unlike transformers, which flatten images into 1D sequences and thus rely heavily on fixed input dimensions, WaveMix processes images in their original 2D structure, making it more flexible for handling images of different sizes and aspect ratios.\\
As for crowd counting, although FFNet and CCTrans achieve competitive performance on the ShanghaiTech Part B (SHB) dataset, several limitations remain, especially when considering deployment on edge devices like Jetson using NVIDIA DeepStream. Both models face challenges in extremely dense scenes due to occlusion and scale variation, which reduce counting accuracy. CCTrans, in particular, requires high computational resources due to its transformer-based architecture, making real-time inference on resource-constrained devices difficult. While FFNet is more lightweight, its performance may be less consistent in complex scenes with varying illumination and crowd distribution. Furthermore, both models rely heavily on training data distribution and may generalize poorly to unseen environments or camera angles in real-world deployment scenarios.\\
As for deepstream, this report detailed a two-stage distributed system designed for real-time people counting, successfully integrating Super Resolution (SR) for enhanced video input quality on a Windows PC with efficient edge inference via NVIDIA DeepStream on a Jetson device. The Windows stage effectively utilized FFmpeg to convert SR-processed images into a consistent RTSP stream, which was then published using Mediamtx. The Jetson device demonstrated its capability to ingest this stream and perform people counting, achieving approximately 25 FPS for a single input source. Performance evaluations further revealed that while the Jetson can handle multiple streams, the per-stream FPS degrades significantly with increasing load, processing four concurrent streams at an average of 7 FPS per stream. This highlights a key trade-off between the number of monitored sources and the achievable real-time processing fidelity on the edge device. The proposed architecture thus provides a flexible framework, though careful consideration of computational resources and desired per-stream performance is crucial for practical deployment, particularly in multi-source scenarios.
\section{Future Work}
Expanding the WaveMix architecture to support other vision tasks like object detection and instance segmentation is a promising area for future work. It would also be beneficial to investigate scaling the model to deeper and wider configurations—by increasing the number of WaveMix blocks. Furthermore, exploring more effective methods to incorporate a wider range of image priors could help lower the computational cost of training neural networks that generalize well across different vision applications.\\
To address the aforementioned limitations, future work can focus on optimizing the models for edge deployment and improving generalization across diverse environments. One potential direction is to apply model compression techniques such as pruning, quantization, or knowledge distillation to reduce the computational overhead of CCTrans while maintaining its accuracy. For FFNet, enhancing robustness through data augmentation or domain adaptation could help improve consistency under varying lighting and crowd distribution conditions. In addition, integrating lightweight attention mechanisms or hybrid CNN-transformer architectures may offer a better balance between performance and efficiency. Finally, expanding the training dataset with real-world footage from deployment scenarios (e.g., Jetson camera streams) could further improve model generalization and adaptability in practice. \\
For examination corrections, we acknowledge the strengths of localization-based approaches in achieving accurate detection results such as APGCC. With the emergence of MAMBA, there is a growing interest in adopting faster and more efficient pyramid-structured models. Inspired by the YOLO-Mamba paper, we are exploring architectures that combine speed with precision, particularly those that leverage hierarchical feature extraction and attention mechanisms for real-time performance.\\
This direction reflects a broader trend toward lightweight yet powerful models that can generalize well across various datasets. Despite the challenges in competing with state-of-the-art (SOTA) models on benchmark datasets, our approach focuses on balancing computational efficiency with detection accuracy. The integration of MAMBA-like mechanisms into localization pipelines opens up promising avenues for both research and real-world applications, especially in scenarios where fast, on-device processing is essential. Further research will involve rigorous evaluation against leading models and fine-tuning architectural components to fully exploit the benefits of MAMBA’s dynamic attention while maintaining the robust localization capabilities inherent in traditional approaches.\\
Our main focus will be on revamping the first part of our system – the part that runs on the Windows PC. Right now, it handles images one by one. We want to upgrade it so it can take in live video from up to four different cameras (as RTSP streams), apply our Super Resolution magic to all of them simultaneously, and then send these enhanced video feeds straight to the Jetson. A big part of this will be cutting out the step where we save images to disk after the SR process. Instead, we'll aim to pipe the high-resolution frames directly from the SR models into the RTSP streaming component. This should speed things up considerably and reduce delays.\\
Think of the Windows PC becoming a central "Super Resolution hub," capable of feeding multiple, crystal-clear video streams to one or more Jetson devices. This will make our system much more practical for real-world setups that often use several cameras.
